% ============================================================
% General Discussion
% ============================================================

\section{General Discussion}

The most consistent finding across all three studies is that frontier
LLMs---particularly GPT-4---are substantially overconfident in their confidence
intervals for quantitative forecasting. GPT-4's ECE ranged from 18.9\%
(Study~3, favorable domain mix) to 43.7\% (Study~1, equity-only), with a
cross-study mean exceeding 35\%. Even Claude, the best-calibrated model in
every study, achieved only 19.7\% ECE in Study~1. This pattern of overconfidence
was robust to domain, horizon, and sample size, confirming it as a systematic
property of these models rather than an artifact of any single design.

These ECE values are comparable to the most severe human expert overconfidence
on record. \citet{BenDavid2013} found that CFOs' 80\% intervals for annual
equity returns achieved only 36--38\% empirical coverage---roughly equivalent
to our Study~1 Claude result (58\% at 80\%). GPT-4's Study~1 performance
(31\% at the 80\% level) is dramatically worse, suggesting that some frontier
LLMs may exhibit super-human overconfidence in certain forecasting contexts.

\paragraph{The horizon effect} Study~2 reveals that prediction horizon
strongly moderates miscalibration, with a divergent model-specific pattern.
The degradation of Claude and GPT-4 is consistent with CI widths anchored
to short-horizon volatility---a failure mode that mirrors human behavior
at longer horizons \citep{Griffin1992}. Gemini's near-perfect calibration
at both 1~day and 22~days suggests its CI widths scale approximately
proportionally to realized volatility at these extremes, though the mechanism
underlying this unusual profile is not yet understood.

\paragraph{The domain-confidence inversion} Study~3's central finding---that
all models overstate certainty in commodity and equity markets while
overstating uncertainty in cryptocurrency and forex---is inconsistent with a
single domain-agnostic miscalibration mechanism. The most plausible
interpretation is that LLMs apply domain-level heuristics inherited from
training corpora: equity research and commodity analysis provide many narrow
price forecasts, encoding an implicit prior that these markets are predictable,
while cryptocurrency and forex commentary may emphasize unpredictability.
The uniform $\mu = 0.25$ on commodities across all three models supports
a shared training-data origin rather than a model-specific artifact.

\paragraph{Model differences} GPT-4 was consistently the most overconfident
model. Its Study~1 ECE (43.7\%) is nearly double Claude's (19.7\%) in the
same task, and its Study~2 ECE at 22~days (49.9\%) approaches the theoretical
maximum for three-level calibration. Claude performed best in every study at
the aggregate level but exhibits severe directional miscalibration in
cryptocurrency ($\mu = 14.4$)---the largest positive outlier in the
meta-knowledge analysis. This shows that even the best-calibrated model
can be dramatically miscalibrated in specific domains. GPT-4's forex failure
(ECE\,=\,44.2\%) is the worst domain-model result in Study~3 and reflects a
compound failure of point-estimate anchoring and CI-width miscalibration that
the MEAD/MAD ratio alone cannot fully capture; both metrics are required for
complete diagnosis.

\paragraph{Limitations} Several limitations constrain our conclusions.
First, Studies~1 and~3 were each conducted on a single day, and Study~3's
domain-specific results may partly reflect idiosyncratic market conditions
on March~25, 2026. Study~2's backtest design mitigates this concern for
equities but introduces potential look-ahead bias if models' training data
extends to the backtest period. Second, statistical precision is limited in
small Study~3 domains ($n = 5$ for commodities, $n = 5$ for crypto,
$n = 8$ for forex). Third, the Claude model version differed between
Studies~1--2 (claude-sonnet-4-20250514) and Study~3 (claude-sonnet-4-6);
cross-study comparisons for Claude should be interpreted with this caveat.
Fourth, NBA outcomes were excluded entirely from Study~3 due to ESPN API
resolution failures. Fifth, models were evaluated without chain-of-thought
or explicit calibration prompting; whether such interventions produce genuine
calibration improvement is an important open question.

\paragraph{Future directions} Future work should: (i)~collect predictions
daily over extended periods to separate stable calibration properties from
single-day conditions; (ii)~extend the domain set to include earnings
forecasts, macroeconomic indicators, and prediction market prices;
(iii)~evaluate calibration-specific prompting and post-hoc recalibration
(conformal prediction, temperature scaling); and (iv)~use larger $n$ within
Study~3 domains to enable robust domain-level inference.


% -----------------------------------------------------------------------
% Conclusion
% -----------------------------------------------------------------------
\section{Conclusion}

We report three preregistered studies examining whether frontier LLMs exhibit
systematic overconfidence in quantitative interval estimates, and how
miscalibration varies with prediction horizon and domain. All three models
were substantially overconfident across all three studies, with GPT-4
consistently the worst (cross-study ECE 19--44\%) and Claude the best.
Study~2 reveals that Claude and GPT-4 become increasingly miscalibrated as
horizon grows from 1 to 22~days, while Gemini maintains near-target coverage
at both ends of the tested range. Study~3 reveals a domain-confidence
inversion: all models overstate certainty in commodity and equity markets
while overstating uncertainty in cryptocurrency and forex---a pattern
consistent with domain-specific miscalibration inherited from training corpora.

These results have direct practical implications. Users who elicit confidence
intervals from LLMs for financial decision support should not assume that
stated uncertainty reflects calibrated epistemic states. The pattern of
miscalibration is most severe in precisely the domains (commodities, equities)
and at precisely the horizons (multi-week) where overconfident forecasting
carries the most severe consequences. Global calibration benchmarks do not
transfer across domains or horizons; domain-stratified and
horizon-stratified evaluation is necessary to characterize model-specific
calibration risk. More broadly, this study demonstrates that the
domain-specificity and horizon-sensitivity of human overconfidence---developed
over four decades of research on human judges---apply with striking regularity
to large language models trained on those same human outputs.


% -----------------------------------------------------------------------
% Acknowledgements
% -----------------------------------------------------------------------
% Acknowledgements moved to main.tex \begin{acknowledgments} block (PNAS format)
